\documentclass[lettersize,journal]{IEEEtran}

\usepackage{amsmath,amsfonts,amssymb}

\usepackage{algorithmic}

\usepackage{algorithm}

\usepackage{array}

\usepackage[caption=false,font=normalsize,labelfont=sf,textfont=sf]{subfig}

\usepackage{textcomp}

\usepackage{stfloats}

\usepackage{url}

\usepackage{verbatim}

\usepackage{graphicx}

\usepackage{cite}

\usepackage{booktabs}

\usepackage{multirow}

\usepackage{bm}

\usepackage{color}

\usepackage{float}

\usepackage{flafter}

\usepackage{makecell}

\usepackage{setspace}

\usepackage[hidelinks,unicode]{hyperref}
\renewcommand{\algorithmicrequire}{\textbf{Input:}}
\renewcommand{\algorithmicensure}{\textbf{Output:}}

\hyphenation{op-tical net-works semi-conduc-tor IEEE-Xplore}

\begin{document}

\title{Multi-Goal Traversal Planning in Cluttered Environments via Improved A* and MMAS-VND-CL}

\author{Author Name~\IEEEmembership{Member,~IEEE}%
\thanks{Author affiliation and correspondence information to be supplied.}%
\thanks{Manuscript received 2026; revised 2026.}}

\markboth{Journal Preprint}%
{Multi-Goal Traversal Planning via Improved A* and MMAS-VND-CL}

\maketitle

\begin{abstract}

A two-layer multi-goal traversal planning framework is proposed to jointly address obstacle-avoiding point-to-point path finding and target visitation ordering in cluttered environments. An improved A* planner based on Jump Point Search (JPS) computes pairwise paths; continuous trajectories are then generated through safety-distance-aware path simplification and arc-length parameterized cubic spline smoothing, and the lengths of the processed trajectories form an obstacle-constrained cost matrix. An MMAS-VND-CL solver that integrates Max--Min Ant System (MMAS), candidate-list-accelerated Variable Neighborhood Descent (VND), pseudo-random transition, and progressive scheduling is subsequently employed to solve the open traveling salesman problem with fixed start and goal via a virtual-node transformation. Experiments show that, on Map 2, the improved A* attains a planning time of 11.85\,ms, lower than classical A* (29.19\,ms) and Dijkstra (202.29\,ms); the processed trajectories are approximately 4.7\%--5.3\% shorter than the raw grid paths of the compared baselines. On the standard closed-loop TSP benchmark kroB100, the solver reaches the known optimal tour cost of 22141 in 97.5\% of 40 runs, with a mean cost about 9.8\% lower than the compared ACO implementation. Ablation studies indicate that VND and the MMAS pheromone bounds are the principal components affecting solution quality and stability. A satellite-image grid scenario further demonstrates the practicality of the framework.

\end{abstract}

\begin{IEEEkeywords}

multi-goal traversal, path planning, ant colony optimization, open traveling salesman problem, obstacle avoidance

\end{IEEEkeywords}

\section{Introduction}

Multi-goal traversal path planning is a fundamental problem in autonomous mobile robotics, with broad applications in autonomous shuttle services, logistics and warehouse management, infrastructure inspection, and search-and-rescue operations \textcolor{blue}{\cite{ref1,ref2,ref33,ref37,ref38}}. In tasks such as passenger pickup and drop-off, equipment inspection, and item delivery, a robot must depart from a designated start location, visit a sequence of intermediate target points, and finally reach a designated goal. The planner must address two subproblems that are distinct yet coupled:

\begin{enumerate}

\item finding collision-free paths between consecutive point pairs in cluttered environments;

\item determining the optimal visitation order of the intermediate target points.

\end{enumerate}

The first subproblem is geometric: it requires reasoning about obstacle geometry and path continuity in physical space. The second is a combinatorial optimization problem, analogous to classical NP-hard problems such as the traveling salesman problem (TSP). The two subproblems interact: the traversal cost between any two points depends on the obstacles lying between them, while the optimal visitation order in turn depends on these pairwise traversal costs.

In general, existing multi-goal navigation methods tend to handle only one of the two subproblems in isolation. A large body of TSP research assumes that pairwise costs are given by Euclidean distance or a known distance matrix \textcolor{blue}{\cite{ref4,ref36}}, overlooking the fact that, in obstacle-rich environments, the shortest collision-free path between two points can be far longer than their straight-line distance \textcolor{blue}{\cite{ref39}}. On the other hand, mature path-planning algorithms such as A*, Dijkstra, and rapidly-exploring random trees (RRT) and their variants can efficiently compute optimal or near-optimal collision-free paths between point pairs on grid maps, but these planners do not optimize the visitation order of multiple goals \textcolor{blue}{\cite{ref5,ref6,ref7}}. A few studies have explored obstacle-aware TSP variants under specific geometric assumptions (e.g., TSP with obstacles, Dubins TSP); however, a general framework that integrates a grid-based global planner with a high-performance TSP solver for multi-goal traversal in cluttered environments is still lacking.

To address these constraints and fully exploit the strengths of both algorithm families, this paper proposes a collaborative optimization framework for multi-goal traversal path planning in cluttered environments. The framework consists of two algorithmic layers connected by a precomputed cost matrix. The first layer is an improved A* global planner based on Jump Point Search (JPS), which exploits grid symmetry to skip redundant node expansions and expands only jump points, thereby greatly reducing the search-tree size; this planner is combined with a safety-aware path simplification procedure and an arc-length parameterized cubic spline smoothing pipeline. The second layer is an improved ant colony optimization (ACO) solver for the open TSP. It converts the open TSP into a standard closed-loop TSP by means of a virtual (dummy) node, thus fully exploiting closed-loop TSP search operators. The solver integrates Max--Min Ant System (MMAS) pheromone management, an Ant Colony System (ACS)-style pseudo-random transition rule, a Variable Neighborhood Descent (VND) local search engine accelerated by a $k$-nearest-neighbor candidate list, and a sigmoid-based progressive local-search schedule with an adaptive stopping criterion.

To tackle the above issues, this paper develops a two-layer planning framework comprising pairwise path search and visitation-order optimization. The former generates feasible paths according to the obstacle layout and computes the corresponding costs; the latter takes this cost matrix as input to determine the visitation order of the targets. On this basis, the path-processing pipeline and the ant-colony ordering solver are improved, respectively. The main contributions are as follows.

\begin{enumerate}

\item A hierarchical solution procedure for multi-goal traversal with fixed start and goal. The lengths of continuous trajectories of obstacle-constrained pairwise paths serve as traversal costs between targets, yielding a cost matrix that links geometric path planning and visitation-order optimization and thereby mitigating the neglect of obstacle-detour costs incurred by Euclidean-distance-based ordering.

\item An improved A* path-planning and post-processing pipeline based on JPS. Incorporating tie-breaking, search-failure fallback, safety-distance-aware path simplification, and arc-length parameterized cubic spline smoothing, the pipeline reduces node expansions while producing continuous paths for cost computation; its planning efficiency and path length are analyzed through grid-map experiments.

\item An MMAS-VND-CL visitation-order solver. A virtual node handles the fixed-endpoint constraint of the open TSP; MMAS pheromone bounds, ACS-style pseudo-random transition, candidate-list-accelerated VND, sigmoid scheduling, and an adaptive stopping criterion are integrated. Parameter-sensitivity and ablation experiments analyze the effects of individual components on solution quality, stability, and runtime.

\end{enumerate}

The remainder of this paper is organized as follows. Section~\textcolor{blue}{\ref{sec:related}} reviews related work on global path planning, TSP metaheuristics, and multi-goal traversal with obstacles; Section~\textcolor{blue}{\ref{sec:form}} formalizes the problem; Section~\textcolor{blue}{\ref{sec:method}} details the methodology, including the improved A* planner, the path post-processing pipeline, the improved ACO solver, and the collaborative optimization framework; Section~\textcolor{blue}{\ref{sec:exp}} reports experimental results along four evaluation dimensions: global path-planning comparison, parameter sensitivity analysis, ablation studies, and multi-goal traversal comparison among different TSP solvers; Section~\textcolor{blue}{\ref{sec:concl}} concludes the paper and outlines future work.

